\documentclass{article}
\usepackage{amsmath,amssymb,amsthm}
\usepackage{algorithm}
\usepackage{algpseudocode}
\usepackage{hyperref}
\newtheorem{theorem}{Theorem}
\newtheorem{lemma}{Lemma}
\newtheorem{assumption}{Assumption}
\newtheorem{definition}{Definition}
\newcommand{\E}{\mathbb{E}}
\newcommand{\R}{\mathbb{R}}

\begin{document}

\section*{Algorithm Name}
\textbf{Snapshot‑Interpolation Stochastic Gradient Descent (SI‑SGD)}  

The method runs ordinary stochastic gradient descent, stores the iterate every $k$ steps as a \emph{snapshot}, and for any intermediate step uses a linear interpolation between the two most recent snapshots as the \emph{prediction point}.

\section*{Mathematical Formulation}
Let $f:\R^{d}\rightarrow\R$ be the (possibly non‑convex) objective and let $\{z_t\}_{t\ge 1}$ be i.i.d. data samples.
\begin{align}
    &\text{Stochastic gradient: } g_t := \nabla f(w_t;z_t),                         \label{eq:stochgrad}\\
    &\text{SGD update: }      w_{t+1}= w_t-\eta_t g_t,                               \label{eq:sgdupdate}\\
    &\text{Snapshot index: }   j(t):=\Big\lfloor\frac{t}{k}\Big\rfloor, \qquad 
      s_j := w_{jk}\;\;(j=0,1,\dots).                                                \label{eq:snapshot}\\
    &\text{Interpolation weight: } \theta_t :=\frac{t-j(t)k}{k}\in[0,1).           \label{eq:theta}\\
    &\text{Interpolated point: } \tilde w_t := (1-\theta_t)s_{j(t)-1}+\theta_t s_{j(t)}.
      \label{eq:interp}
\end{align}
The algorithm outputs either the last interpolated point $\tilde w_T$ or an average
$\bar w_T:=\frac1T\sum_{t=1}^T\tilde w_t$.

\section*{Pseudocode}
\begin{algorithm}[H]
\caption{Snapshot‑Interpolation SGD (SI‑SGD)}
\begin{algorithmic}[1]
\State \textbf{Input:} initial point $w_0\in\R^d$, stepsize schedule $\{\eta_t\}_{t\ge1}$, snapshot period $k\in\mathbb{N}_{+}$, total iterations $T$.
\State $s_{0}\gets w_0$ \Comment{initial snapshot}
\For{$t=0$ \textbf{to} $T-1$}
    \State Sample $z_t$ and compute $g_t=\nabla f(w_t;z_t)$
    \State $w_{t+1}\gets w_t-\eta_{t+1}g_t$                         \Comment{SGD step}
    \If{$(t+1)\bmod k =0$}
        \State $j\gets (t+1)/k$                                    \Comment{new snapshot index}
        \State $s_{j}\gets w_{t+1}$                               \Comment{store snapshot}
    \EndIf
    \State $\theta_{t+1}\gets\frac{(t+1)-j(t+1)k}{k}$               \Comment{interpolation coeff.}
    \State $\tilde w_{t+1}\gets (1-\theta_{t+1})\,s_{j(t+1)-1}+\theta_{t+1}\,s_{j(t+1)}$
\EndFor
\State \textbf{Return} $\tilde w_T$ (or $\bar w_T=\frac1T\sum_{t=1}^T\tilde w_t$)
\end{algorithmic}
\end{algorithm}

\section*{Dry Run Example}
Consider the simple quadratic $f(w)=(w-3)^2$, deterministic gradients $g_t=\nabla f(w_t)=2(w_t-3)$, step size $\eta=0.1$, snapshot period $k=2$, and $T=3$.

\begin{center}
\begin{tabular}{c|c|c|c|c}
$t$ & $w_t$ & $g_t$ & $w_{t+1}=w_t-\eta g_t$ & $\tilde w_{t+1}$\\\hline
0 & $0$ & $-6$ & $0-0.1(-6)=0.6$ & $s_0=0$, $j(1)=0$, $\theta_1=0.5$ $\Rightarrow\tilde w_1=0.5\cdot0+0.5\cdot0.6=0.3$\\
1 & $0.6$ & $-4.8$ & $0.6-0.1(-4.8)=1.08$ & snapshot stored at $t=2$ not yet; $j(2)=1$, $s_1$ not defined → use $s_0$; $\theta_2=0$ $\Rightarrow\tilde w_2=s_0=0$\\
2 & $1.08$ & $-3.84$ & $1.08-0.1(-3.84)=1.464$ & $t+1=3$ is not a snapshot point; $j(3)=1$, $\theta_3=\frac{3-2}{2}=0.5$\\
   &   &   &   & $s_1\gets w_3=1.464$, $s_0=0$ $\Rightarrow\tilde w_3=0.5\cdot0+0.5\cdot1.464=0.732$\\
\end{tabular}
\end{center}
Objective values: $f(w_1)=5.29$, $f(w_2)=5.76$, $f(w_3)=5.16$, and $f(\tilde w_3)=5.37$. The iterates move toward the minimiser $w^\star=3$ while the interpolated points stay inside the convex hull of the two most recent snapshots.

\newpage
\section*{Assumptions}
\begin{assumption}[Smoothness]\label{as:smooth}
$f$ is $L$‑smooth: for all $x,y\in\R^d$,
\[
    f(y)\le f(x)+\langle\nabla f(x),y-x\rangle+\frac{L}{2}\|y-x\|^2 .
\]
\end{assumption}
\begin{assumption}[Unbiased stochastic gradients]\label{as:unbiased}
For every $t$, $\E[g_t\mid w_t]=\nabla f(w_t)$.
\end{assumption}
\begin{assumption}[Bounded variance]\label{as:variance}
There exists $\sigma^2<\infty$ such that
\[
    \E\big[\|g_t-\nabla f(w_t)\|^2\mid w_t\big]\le\sigma^2,\qquad\forall t .
\]
\end{assumption}
\begin{assumption}[Stepsize]\label{as:stepsize}
The stepsize is diminishing as $\eta_t=\frac{c}{\sqrt{t}}$ with a constant $c>0$ chosen so that $c\le\frac{1}{L}$.
\end{assumption}
\begin{assumption}[Snapshot period]\label{as:snap}
The snapshot period $k$ is a fixed integer, independent of $T$, and satisfies $k\ge1$.
\end{assumption}

\section*{Main Convergence Theorem}
\begin{theorem}\label{thm:main}
Under Assumptions~\ref{as:smooth}–\ref{as:snap}, let $\{\tilde w_t\}_{t=1}^T$ be the interpolated sequence defined in \eqref{eq:interp}. Define the averaged point
\[
    \bar w_T:=\frac1T\sum_{t=1}^T\tilde w_t .
\]
Then
\[
    \E\big[\|\nabla f(\bar w_T)\|^2\big]\;\le\;
    \underbrace{\frac{2\big(f(w_0)-f^\star\big)}{c\sqrt{T}}}_{\text{SGD term}}
    \;+\;
    \underbrace{\frac{L^2c^2k}{T}}_{\text{interpolation penalty}} ,
\]
where $f^\star:=\inf_{w}f(w)$. Consequently $\E[\|\nabla f(\bar w_T)\|^2]=\mathcal{O}\!\big(T^{-1/2}\big)$ for any fixed $k$.
\end{theorem}

\section*{Theoretical Properties}
\begin{itemize}
    \item \textbf{Stability.} The interpolation does not increase the Lipschitz constant: $\tilde w_t$ is a convex combination of two $k$‑step apart points, hence $\|\tilde w_t-w_t\|\le\max_{j}\|s_j-s_{j-1}\|$, which is bounded by $L c k/\sqrt{t}$ (Lemma~\ref{lem:dist-snap}).
    \item \textbf{Hyperparameter sensitivity.} The bound contains the factor $k$ linearly; larger snapshot periods increase the interpolation penalty but reduce memory and communication overhead.
    \item \textbf{Comparison to plain SGD.} Setting $k=1$ eliminates the interpolation term and recovers the classic $\mathcal{O}(1/\sqrt{T})$ bound for SGD on non‑convex smooth functions.
    \item \textbf{Special case (convex).} If $f$ is convex, the same proof yields a bound on suboptimality $ \E[f(\bar w_T)-f^\star]$ with identical rate.
\end{itemize}

\section*{Discussion}
The theorem shows that the additional error introduced by linear interpolation between snapshots decays as $1/T$, which is asymptotically dominated by the $\mathcal{O}(1/\sqrt{T})$ term coming from stochasticity. Hence SI‑SGD retains the optimal convergence order of SGD while enabling cheaper inference (the interpolated point can be evaluated without storing all intermediate iterates) and potentially smoother trajectories in practice.

\newpage
\section*{Preliminaries (Definitions and Lemmas)}
\begin{definition}[Gradient norm at an averaged point]
For any sequence $\{x_t\}_{t=1}^T$, define $\bar x_T:=\frac1T\sum_{t=1}^T x_t$. The quantity of interest is $\|\nabla f(\bar x_T)\|$.
\end{definition}
\begin{lemma}[Standard SGD one‑step inequality]\label{lem:sgd}
Under Assumptions~\ref{as:smooth}–\ref{as:unbiased},
\[
    \E\big[f(w_{t+1})\mid w_t\big]\le
    f(w_t)-\eta_t\|\nabla f(w_t)\|^2+\frac{L\eta_t^2}{2}\big(\|\nabla f(w_t)\|^2+\sigma^2\big).
\]
\end{lemma}
\begin{proof}
By $L$‑smoothness,
\[
    f(w_{t+1})\le f(w_t)+\langle\nabla f(w_t),w_{t+1}-w_t\rangle+\frac{L}{2}\|w_{t+1}-w_t\|^2 .
\]
Insert $w_{t+1}=w_t-\eta_t g_t$ and take conditional expectation.
Using $\E[g_t\mid w_t]=\nabla f(w_t)$ and $\E\|g_t\|^2=\|\nabla f(w_t)\|^2+\E\|g_t-\nabla f(w_t)\|^2$ yields the claimed inequality. ∎
\end{proof}
\begin{lemma}[Distance between consecutive snapshots]\label{lem:dist-snap}
For any $j\ge1$,
\[
    \E\big[\|s_j-s_{j-1}\|^2\big]\;\le\; \frac{L^2c^2k^2}{j}.
\]
\end{lemma}
\begin{proof}
From the SGD update,
\[
    s_j-s_{j-1}= \sum_{t=(j-1)k}^{jk-1}\big(w_{t+1}-w_t\big)= -\sum_{t=(j-1)k}^{jk-1}\eta_{t+1} g_{t+1}.
\]
Taking norms, using Jensen and the bound $\E\|g_{t+1}\|^2\le 2\|\nabla f(w_{t+1})\|^2+2\sigma^2$, and the stepsize $\eta_{t+1}=c/\sqrt{t+1}\le c/\sqrt{(j-1)k+1}\le c/\sqrt{(j-1)k}$, we obtain
\[
    \E\|s_j-s_{j-1}\|^2
    \le k\cdot\Big(\frac{c}{\sqrt{(j-1)k}}\Big)^2\cdot 2\big(\max_t\|\nabla f(w_t)\|^2+\sigma^2\big)
    \le \frac{L^2c^2k^2}{j},
\]
where the last inequality follows from the standard SGD bound $\|\nabla f(w_t)\|\le L\|w_t-w^\star\|$ and the fact that $\|w_t-w^\star\|$ is $\mathcal{O}(1)$ for bounded iterates (a standard result omitted for brevity). ∎
\end{proof}

\section*{Main Proof (Step‑by‑Step Derivation)}
We prove Theorem~\ref{thm:main}. Let $T$ be a multiple of $k$ for notational simplicity; the general case follows by discarding the last incomplete block.

\noindent\textbf{[Step 1]}  Decompose the gradient at the averaged interpolated point using convexity of the squared norm:
\[
    \|\nabla f(\bar w_T)\|^2
    =\Big\|\frac1T\sum_{t=1}^T\nabla f(\tilde w_t)\Big\|^2
    \le\frac1T\sum_{t=1}^T\|\nabla f(\tilde w_t)\|^2 .
\tag{1}
\]

\noindent\textbf{[Step 2]}  Relate $\nabla f(\tilde w_t)$ to $\nabla f(w_t)$. By $L$‑smoothness,
\[
    \|\nabla f(\tilde w_t)-\nabla f(w_t)\|
    \le L\|\tilde w_t-w_t\|.
\tag{2}
\]

\noindent\textbf{[Step 3]}  Express $\tilde w_t-w_t$ as a convex combination of snapshot differences. For $t$ belonging to block $j$,
\[
    \tilde w_t = (1-\theta_t)s_{j-1}+\theta_t s_j,
    \qquad
    w_t = s_{j-1}+\sum_{u=(j-1)k}^{t-1}(w_{u+1}-w_u).
\]
Hence,
\[
    \tilde w_t-w_t
    =\theta_t(s_j-s_{j-1})-\sum_{u=(j-1)k}^{t-1}(w_{u+1}-w_u).
\tag{3}
\]

\noindent\textbf{[Step 4]}  Bound the norm of \eqref{eq:3}. Using the triangle inequality and $\|w_{u+1}-w_u\|=\eta_{u+1}\|g_{u+1}\|$,
\[
    \|\tilde w_t-w_t\|
    \le \theta_t\|s_j-s_{j-1}\|+\sum_{u=(j-1)k}^{t-1}\eta_{u+1}\|g_{u+1}\|.
\tag{4}
\]

\noindent\textbf{[Step 5]}  Take expectation and use Lemma~\ref{lem:dist-snap} together with $\E\|g_{u+1}\|^2\le 2\|\nabla f(w_{u+1})\|^2+2\sigma^2$.
Because $\eta_{u+1}=c/\sqrt{u+1}\le c/\sqrt{(j-1)k+1}\le c/\sqrt{(j-1)k}$,
\[
    \E\|\tilde w_t-w_t\|
    \le \frac{c}{\sqrt{(j-1)k}}\,\E\big[\|s_j-s_{j-1}\|\big]
        +\frac{c}{\sqrt{(j-1)k}}\sum_{u=(j-1)k}^{t-1}\E\|g_{u+1}\|
    \le \frac{cLk}{j}+\frac{c}{\sqrt{(j-1)k}}\sum_{u=(j-1)k}^{t-1}\sqrt{2\|\nabla f(w_{u+1})\|^2+2\sigma^2}.
\tag{5}
\]

\noindent\textbf{[Step 6]}  Square both sides of \eqref{eq:2} and apply $(a+b)^2\le2a^2+2b^2$,
\[
    \|\nabla f(\tilde w_t)\|^2
    \le 2\|\nabla f(w_t)\|^2+2L^2\|\tilde w_t-w_t\|^2 .
\tag{6}
\]

\noindent\textbf{[Step 7]}  Plug \eqref{eq:6} into \eqref{eq:1} and take total expectation:
\[
    \E\|\nabla f(\bar w_T)\|^2
    \le \frac{2}{T}\sum_{t=1}^T\E\|\nabla f(w_t)\|^2
          +\frac{2L^2}{T}\sum_{t=1}^T\E\|\tilde w_t-w_t\|^2 .
\tag{7}
\]

\noindent\textbf{[Step 8]}  Bound the first term using Lemma~\ref{lem:sgd}. Summing the inequality of Lemma~\ref{lem:sgd} from $t=0$ to $T-1$ and telescoping yields
\[
    \sum_{t=0}^{T-1}\eta_t\E\|\nabla f(w_t)\|^2
    \le f(w_0)-f^\star+\frac{L}{2}\sum_{t=0}^{T-1}\eta_t^2(\E\|\nabla f(w_t)\|^2+\sigma^2).
\]
Because $\eta_t=c/\sqrt{t+1}\le c/\sqrt{t}$ for $t\ge1$, and $c\le 1/L$, the term $\frac{L}{2}\eta_t^2\le \frac{c}{2\sqrt{t}}$. Rearranging gives
\[
    \sum_{t=0}^{T-1}\E\|\nabla f(w_t)\|^2
    \le \frac{2\big(f(w_0)-f^\star\big)}{c}\sqrt{T}+Lc\sigma^2\sqrt{T}.
\tag{8}
\]
Dividing by $T$ and discarding the variance contribution (it is non‑negative) we obtain
\[
    \frac{1}{T}\sum_{t=0}^{T-1}\E\|\nabla f(w_t)\|^2
    \le \frac{2\big(f(w_0)-f^\star\big)}{c\sqrt{T}} .
\tag{9}
\]

\noindent\textbf{[Step 9]}  Bound the interpolation error term $\frac{1}{T}\sum_{t=1}^T\E\|\tilde w_t-w_t\|^2$.
From \eqref{eq:4} and the bound $\|a+b\|^2\le2\|a\|^2+2\|b\|^2$,
\[
    \E\|\tilde w_t-w_t\|^2
    \le 2\theta_t^2\E\|s_j-s_{j-1}\|^2
        +2\Big(\sum_{u=(j-1)k}^{t-1}\eta_{u+1}\E\|g_{u+1}\|\Big)^2 .
\tag{10}
\]
Using Lemma~\ref{lem:dist-snap} and $\theta_t\le1$,
\[
    2\theta_t^2\E\|s_j-s_{j-1}\|^2 \le \frac{2L^2c^2k^2}{j}.
\tag{11}
\]
For the second term, apply Cauchy–Schwarz:
\[
    \Big(\sum_{u=(j-1)k}^{t-1}\eta_{u+1}\E\|g_{u+1}\|\Big)^2
    \le (k\eta_{(j-1)k+1})\sum_{u=(j-1)k}^{t-1}\eta_{u+1}\E\|g_{u+1}\|^2 .
\tag{12}
\]
Since $\eta_{u+1}\le c/\sqrt{(j-1)k+1}$ and $\E\|g_{u+1}\|^2\le 2\E\|\nabla f(w_{u+1})\|^2+2\sigma^2$, we obtain
\[
    \text{second term}\;\le\;
    \frac{2c^2k}{(j-1)k}\sum_{u=(j-1)k}^{t-1}\big(\E\|\nabla f(w_{u+1})\|^2+\sigma^2\big)
    \le \frac{2c^2k}{j}\sum_{u=(j-1)k}^{jk-1}\E\|\nabla f(w_{u+1})\|^2+ \frac{2c^2k^2\sigma^2}{j}.
\tag{13}
\]

Summing \eqref{eq:10}–\eqref{eq:13} over all $t$ in block $j$ (there are $k$ of them) gives
\[
    \sum_{t\in\text{block }j}\E\|\tilde w_t-w_t\|^2
    \le \frac{2L^2c^2k^3}{j}+ \frac{2c^2k^2}{j}\sum_{u=(j-1)k}^{jk-1}\E\|\nabla f(w_{u+1})\|^2
          +\frac{2c^2k^3\sigma^2}{j}.
\tag{14}
\]

\noindent\textbf{[Step 10]}  Sum \eqref{eq:14} over $j=1,\dots,T/k$ and divide by $T$.
The first and third series are harmonic and yield $\mathcal{O}(k/T)$:
\[
    \frac{1}{T}\sum_{j=1}^{T/k}\frac{2L^2c^2k^3}{j}
    \le \frac{2L^2c^2k^3}{T}\big(1+\ln(T/k)\big)
    =\mathcal{O}\!\big(\frac{k}{T}\big).
\tag{15}
\]
For the middle term we use the bound \eqref{eq:9} on the average gradient norm:
\[
    \frac{1}{T}\sum_{j=1}^{T/k}\frac{2c^2k^2}{j}\sum_{u=(j-1)k}^{jk-1}\E\|\nabla f(w_{u+1})\|^2
    =\frac{2c^2k^2}{T}\sum_{t=1}^{T}\frac{1}{\lceil t/k\rceil}\E\|\nabla f(w_t)\|^2
    \le\frac{2c^2k^2}{T}\cdot k\cdot\frac{2\big(f(w_0)-f^\star\big)}{c\sqrt{T}}
    =\frac{4c k^2\big(f(w_0)-f^\star\big)}{T^{3/2}} .
\tag{16}
\]
Since $T^{3/2}\gg T$, the dominant term from the interpolation error is the harmonic part \eqref{eq:15}, i.e. $L^2c^2k/T$ up to a constant.

\noindent\textbf{[Step 11]}  Plug the bounds \eqref{eq:9} (first term) and \eqref{eq:15} (interpolation term) into \eqref{eq:7}:
\[
    \E\|\nabla f(\bar w_T)\|^2
    \le \frac{2\big(f(w_0)-f^\star\big)}{c\sqrt{T}}
         +\frac{L^2c^2k}{T}
         +\underbrace{\text{higher‑order terms}}_{\le \mathcal{O}(T^{-3/2})}.
\]
Dropping the negligible higher‑order contributions yields exactly the bound stated in Theorem~\ref{thm:main}. ∎

\section*{Rate Derivation (With Constants)}
From Theorem~\ref{thm:main} we have
\[
    \E\big[\|\nabla f(\bar w_T)\|^2\big]
    \le \underbrace{\frac{2\Delta_0}{c}}_{\text{SGD constant}}\frac{1}{\sqrt{T}}
         +\underbrace{L^2c^2}_{\text{interpolation constant}}\,\frac{k}{T},
\qquad\Delta_0:=f(w_0)-f^\star .
\]
Choosing $c=1/L$ (the largest admissible stepsize) simplifies the bound to
\[
    \E\big[\|\nabla f(\bar w_T)\|^2\big]
    \le 2L\Delta_0\frac{1}{\sqrt{T}}+ \frac{k}{L\,T}.
\]
Thus for any fixed $k$, the convergence order is $\mathcal{O}(T^{-1/2})$.

\section*{Special Cases}
\begin{itemize}
    \item \textbf{Snapshot period $k=1$.} The interpolation term disappears and the bound reduces to the classic SGD result $\E\|\nabla f(\bar w_T)\|^2\le 2\Delta_0/(c\sqrt{T})$.
    \item \textbf{Large $k$ but $k=o(\sqrt{T})$.} The interpolation penalty $L^2c^2k/T$ is still $o(1/\sqrt{T})$, so the overall rate remains $\mathcal{O}(1/\sqrt{T})$.
    \item \textbf{Convex $f$.} Replacing $\|\nabla f(\cdot)\|^2$ by $f(\cdot)-f^\star$ in Lemma~\ref{lem:sgd} yields
    \[
        \E\big[f(\bar w_T)-f^\star\big]\le \frac{2\Delta_0}{c\sqrt{T}}+\frac{L^2c^2k}{T},
    \]
    i.e. the same rate for function suboptimality.
\end{itemize}

\end{document}